\paragraph{}
\hspace-11ptSäkerheten är alltid av högsta prioritering för en lokförare; hen måste upprätthålla ett koncentrerat fokus framåt under körning.
Förutom att fokusera på banan framåt måste även tågets framförande (pådrag, rullning, broms etc.), status på ombordsystem, diverse inkommande och utgående samtal och meddelanden (från järnvägsföretagets trafikledning, Trafikverkets trafikledning samt tjänstgörande ombordpersonal) skänkas fokus och det är upp till föraren att prioritera och göra rätt beroende på vad situationen kräver.
När då ytterligare ett moment som har hög prioritet tillförs bland nämnda distraktioner bör det ske på ett sådant sätt att föraren \emph{vill} och \emph{kan} följa de anvisningar som förordas via C-DAS\@.
Trafikverkets pilottester har påvisat att förarens förståelse av kontexten för den nya angivna körplanen kan spela in i viljan att realisera den~\cite{Angel2020, Lindmark2022}.
I förlängningen kan detta påverka järnvägsföretagens implementeringstakt av C-DAS\@.
%Att framföra tåg är ett hantverk; e
En skicklig förare har efter flera års yrkesutövande djuplodande kännedom om fordonet som framförs och banans särskilda egenskaper och därmed vilken körning som fordras därefter [källa behövs].
En förare som förstår \emph{varför} inlärd erfarenhetsbaserad hastighetsprofil kan behöva frångås har större vilja att köra inom den mottagna anvisningen.
Då realtidsjusteringarna av tidtabeller ofta berör flera tåg kommer de som inte tar emot och följer dessa anvisningar få en stor negativ inverkan på trafikflödet~\cite{Joborn2023}.
Dels för att det inte går att se hur tåget framförs men även att det inte kan delges ny tidtabellsinformation.
Det ligger därför i Trafikverkets intresse att så många järnvägsföretag som möjligt implementerar en C-DAS-lösning.
